\section{Evaluación automática del habla}
\textcolor{teal}{Métodos de evaluación de pronunciación, fluidez y prosodia mediante técnicas de procesamiento del habla. Revisión de los artículos académicos y los sistemas existentes con el análisis y evaluación de las características del habla.}
\textcolor{teal}{
    Qué evaluar en habla:
    \begin{itemize}
        \item pronunciación
        \item fluidez
        \item ritmo (prosodia)
    \end{itemize}
    Métodos clásicos:
    \begin{itemize}
        \item análisis acústico
        \item modelos fonéticos
    \end{itemize}
    Problemas:
    \begin{itemize}
        \item alta variabilidad
        \item dependencia del idioma
        \item dificultad con niños (clave si lo quieres mencionar indirectamente)
    \end{itemize}
    Conexión con LLMs:
    \begin{itemize}
        \item LLM no trabaja directamente con audio
        \item depende de ASR
    \end{itemize}
    Idea clave -> la evaluación de habla se convierte en evaluación de texto + pérdida de información
}